\section*{Asesmen}
\addcontentsline{toc}{section}{Asesmen}

Asesmen Bab III mencakup Kuis (5\%) dan bagian awal Tugas Teori
(8\%). Fokusnya adalah pemahaman arsitektur dan aliran informasi
antar-fase.

\begin{table}[H]
\centering
\caption{Instrumen asesmen Bab III}
\label{tab:ch03-asesmen}
\begin{tabularx}{\textwidth}{@{}lXl@{}}
\toprule
\textbf{Instrumen} & \textbf{Indikator} & \textbf{Bobot} \\
\midrule
Kuis & Membedakan translator, compiler, interpreter, dan assembler &
5\% \\
Kuis & Menjelaskan fase kompilasi dan peran tabel simbol & 5\% \\
Tugas Teori & Menganalisis trade-off one-pass dan multi-pass & 8\% \\
Praktikum M0 & Repositori dan spesifikasi v0.1 lengkap & prasyarat \\
\bottomrule
\end{tabularx}
\end{table}

Kriteria kelulusan: mahasiswa dapat menjelaskan pipeline kompilasi,
menempatkan tabel simbol dan penanganan kesalahan pada pipeline, serta
memilih strategi kompilasi yang sesuai dengan kebutuhan.
